% Figure: Wien-bridge oscillator. A non-inverting op-amp amplifier of gain 3
% with a series-RC and parallel-RC frequency-selective network feeding the
% positive input from the output. Drawn in plain tikz: op-amp triangle,
% the series and shunt RC arms of the positive-feedback network, and the two
% resistors of the negative-feedback (gain-setting) arm.
\begin{figure}[htbp]
\centering
\begin{tikzpicture}[
  line width=0.8pt,
  every node/.style={font=\small},
  dot/.style={circle, fill=engblack, inner sep=1.3pt},
]
% op-amp triangle
\draw[engblue, line width=1.0pt] (5.0,1.0) -- (5.0,3.0) -- (6.6,2.0) -- cycle;
\node[engblue] at (5.35,2.0) {$A$};
\node[anchor=east, font=\scriptsize] at (5.0,2.6) {$+$};
\node[anchor=east, font=\scriptsize] at (5.0,1.4) {$-$};
% output node
\draw (6.6,2.0) -- (8.2,2.0);
\node[dot] (out) at (7.6,2.0) {};
\node[anchor=south, engblue] at (7.6,2.15) {$v_{\mathrm{out}}$};
% --- positive-feedback frequency-selective network (top path) ---
% from output back to the + input via series RC, with shunt RC to ground.
\draw (7.6,2.0) -- (7.6,3.6) -- (5.9,3.6);
% series arm: R in series with C (series RC)
\draw[engorange, fill=engorange!8] (5.9,3.4) rectangle (5.1,3.8); % series R
\node[engorange, font=\scriptsize] at (5.5,3.6) {$R$};
\draw (5.1,3.6) -- (4.6,3.6);
\draw[line width=1.1pt] (4.6,3.4) -- (4.6,3.8); % series C plate 1
\draw[line width=1.1pt] (4.45,3.4) -- (4.45,3.8); % series C plate 2
\node[font=\scriptsize] at (4.5,4.0) {$C$};
\draw (4.45,3.6) -- (3.6,3.6) -- (3.6,2.6);
% the + input tap node
\node[dot] (ptap) at (3.6,2.6) {};
\draw (3.6,2.6) -- (5.0,2.6);
\node[anchor=south east, font=\scriptsize] at (3.55,2.65) {$v_+$};
% shunt arm: R parallel C from + tap to ground
\draw (3.6,2.6) -- (3.6,1.9);
% shunt R
\draw[engorange, fill=engorange!8] (3.4,1.5) rectangle (3.8,1.9);
\node[engorange, font=\scriptsize] at (3.6,1.7) {$R$};
\draw (3.6,1.5) -- (3.6,1.1);
% shunt C branch (parallel), offset to the left
\draw (3.6,1.9) -- (2.9,1.9) -- (2.9,1.55);
\draw[line width=1.1pt] (2.7,1.55) -- (3.1,1.55);
\draw[line width=1.1pt] (2.7,1.4) -- (3.1,1.4);
\draw (2.9,1.4) -- (2.9,1.1) -- (3.6,1.1);
\node[font=\scriptsize] at (2.55,1.5) {$C$};
% ground for shunt arm
\draw[enggray] (3.6,1.1) -- (3.6,0.6);
\draw[enggray] (3.35,0.6) -- (3.85,0.6);
\draw[enggray] (3.45,0.5) -- (3.75,0.5);
\draw[enggray] (3.53,0.4) -- (3.67,0.4);
% --- negative-feedback gain-setting arm (bottom path) ---
% output to - input via Rf, and - input to ground via Rg
\draw (7.6,2.0) -- (7.6,0.6) -- (6.0,0.6);
\draw[engred, fill=engred!8] (6.0,0.4) rectangle (5.2,0.8); % Rf
\node[engred, font=\scriptsize] at (5.6,0.6) {$R_f$};
\draw (5.2,0.6) -- (4.4,0.6);
\node[dot] (ntap) at (4.4,0.6) {};
% - input connection
\draw (4.4,0.6) -- (4.4,1.4) -- (5.0,1.4);
% Rg to ground
\draw[engred, fill=engred!8] (4.2,0.0) rectangle (4.6,-0.4);
\node[engred, font=\scriptsize] at (4.4,-0.2) {$R_g$};
\draw (4.4,0.6) -- (4.4,0.0);
\draw[enggray] (4.4,-0.4) -- (4.4,-0.7);
\draw[enggray] (4.15,-0.7) -- (4.65,-0.7);
\draw[enggray] (4.25,-0.8) -- (4.55,-0.8);
\end{tikzpicture}
\caption[Wien-bridge oscillator]{The Wien-bridge oscillator: a non-inverting
op-amp amplifier whose positive input is driven from the output through a
frequency-selective network of a series resistor-capacitor arm and a shunt
resistor-capacitor arm. The network passes a maximum fraction of the output,
one third, with zero phase shift at the single frequency where the series and
shunt reactances balance, so the loop meets the Barkhausen condition there and
nowhere else. The negative-feedback arm sets the amplifier gain to just above
three so the loop starts, and a small amplitude-dependent element in that arm
pulls the average gain back to exactly three once the oscillation grows, holding
a clean sinusoid. The Wien bridge is the classic low-distortion audio-frequency
sine source, the counterpart to the LC Colpitts oscillator at radio
frequencies.}
\label{fig:vol04ch09:wien-bridge}
\end{figure}
